\documentclass[10pt]{article}
\usepackage[utf8]{inputenc}
\usepackage[T1]{fontenc}
\usepackage[version=4]{mhchem}
\usepackage{stmaryrd}
\usepackage[margin=0.5in, tmargin=1in]{geometry}
\usepackage{xcolor}
\usepackage{graphicx}
\usepackage{multicol}

\usepackage{fancyhdr}
\pagestyle{fancy}
\fancyhf{}
\fancyhead[L]{Math 1A Spring 2024 Notes}
\fancyhead[R]{Alexis Fernandez}
\renewcommand{\headrulewidth}{0.4pt}

\usepackage{enumitem}
\setlist[enumerate]{leftmargin=*}

\usepackage{amsmath, amsfonts, amssymb}

\author{}
\date{}


\begin{document}

\section*{\centering Note 3 - Transformations of Functions}

\subsection*{Transformations}

\begin{itemize}
    \item \textbf{Four Shifts}
        \begin{multicols}{2}
        \begin{itemize}[label=$\circ$]
            \item Up by $c$ units: $f(x) + c$
            \item Down by $c$ units: $f(x) - c$
            \item Right by $c$ units: $f(x-c)$
            \item Left by $c$ units: $f(x + c)$
        \end{itemize}
        \end{multicols}
        
    \item \textbf{Two Reflections}
        \begin{multicols}{2}
        \begin{itemize}[label=$\circ$]
            \item Through $x$-axis: $-f(x)$
            \item Through $y$-axis: $f(-x)$
        \end{itemize}
        \end{multicols}
    
    \item \textbf{Vertical Stretch/Shrink}
        \begin{itemize}[label=$\circ$]
            \item Stretch by a factor of $c$ units: $cf(x)$ where $c > 1$
            \item Shrink by a factor of $c$ units: $cf(x)$ where $0 < c < 1$ \\
            \textit{Recall: if $c < 0$, then it is a reflection combined with a stretch/shrink!}
        \end{itemize}
    
    \item \textbf{Horizontal Stretch/Shrink}
        \begin{itemize}[label=$\circ$]
            \item Stretch by a factor of $c$ units: $f(cx)$ where $0 < c < 1$
            \item Shrink by a factor of $c$ units: $f(cx)$ where $c > 1$ \\
        \end{itemize}

    \item \textbf{Mother Functions}: Base functions that are often transformed.
        \begin{multicols}{2}
        \begin{itemize}[label=$\circ$]
            \item $y = |x|$
            \item $y = \sqrt{x}$
            \item $y=x^2$
            \item $y=x^3$
        \end{itemize}
        \end{multicols}
    
    
    \item \textbf{Order of Operations}: Given a Mother Function, perform transformations in this order:
    \indent \begin{enumerate}
        \item All stretches or shrinks
        \item All reflections
        \item All shifts
    \end{enumerate}
    

\end{itemize}


\newpage

\subsection*{Composite Functions}

\begin{itemize}
    \item $f$ composed with $g$: $(f \circ g)(x) = f(g(x))$
    \item $g$ composed with $f$: $(g \circ f)(x) = g(f(x))$
    \item $f$ composed with $f$: $(f \circ f)(x) = f(f(x))$
    \item $g$ composed with $g$: $(g \circ g)(x) = g(g(x))$
\end{itemize}

For composite functions, evaluate the inner function 
\textbf{Method 1:} Fractional (Rational) Functions\\
{\color{magenta}
\textbf{Procedure:} Find values of $x$ that make the denominator 0.}
\begin{itemize}
    \item Set denominator equal to 0 and solve for $x$.
    \item The domain is the set of all real numbers except the values of x found in previous step. \\
    \textit{Recall: A denominator of 0 makes the function undefined.}
\end{itemize}
{\color{blue}
\textbf{Example:}} $h(x) = \dfrac{3-x}{2x+1}$

    \begin{align*}
    2x+1 &= 0 \qquad \text{\color{magenta}{Set the denominator equal to 0.}}\\ 
    2x &= -1 \qquad \text{\color{magenta}{Solve for x.}}\\ 
    x &= -\frac{1}{2}
    \end{align*}

    Thus, the domain is everywhere $x \neq -\frac{1}{2}$, since plugging in $-\frac{1}{2}$ would lead to invalid outputs. \\
    \indent In interval notation, this is written as $$(-\infty, -\frac{1}{2})\cup (-\frac{1}{2}, \infty)$$ 
    \begin{quote} 
    \textit{Note: The parentheses are used to show that the interval is not inclusive of $-\frac{1}{2}$ in the domain, but includes numbers greater than $-0.4\overline{999}$ and less than $-0.5000...1$. The $\cup$ represents the union of two subdomains.} \\
    \end{quote} 

\noindent \textbf{Method 2:} Even Root Functions \quad \textit{i.e. $\sqrt{x}, \sqrt[4]{x}, \sqrt[6]{x}$...}\\
{\color{magenta}
\textbf{Procedure:} Find values of $x$ that make the expression inside the root greater than or equal to 0.}
\begin{itemize}
    \item Set the expression inside the radical it greater than or equal to 0, then solve the polynomial inequality.
    \item The solution to the inequality is the domain. \\
    \textit{Recall: Negative numbers inside even roots lead to imaginary solutions.}
\end{itemize}
{\color{blue}
\textbf{Example:}} $g(x) = \sqrt{x^2-x-6}$

    \begin{align*}
    x^2-x-6 &\geq 0 \qquad \text{\color{magenta}{Set the expression greater than or equal to 0.}}\\ 
    (x-3)(x+2) &= 0 \qquad \text{\color{magenta}{Solve for $x$. \textit{Note: the inequality is temporarily changed to an equality.}}}\\ 
    x &= 3, -2 \qquad \text{\color{magenta}{Thus, the inequality is equal to 0 when $x = 3$ or $-2$.}}
    \end{align*}

    \begin{quote} 
     When $x = 3$ or $-2$, then we know that the inequality is equal to 0. However, we must also check whether the inequality is true for the intervals where $x < -2, -2 < x<3$, and $x > 3$. For each interval, plug in a test point for $x$ and see if the inequality holds true.
    \end{quote}
    \begin{itemize}
        \item For $x < -2$: plugging in $x = -3$ gives $(-3)^2+ 3 -6 \geq 0$. When simplified, this gives $ 6 \geq 0$, which is true. Thus, real numbers less than or equal to -2 are in the domain.
        \item For $-2 < x < 3$: plugging in $x = 0$ gives $(0)^2 - 0 -6 \geq 0$. When simplified, this gives $ -6 \geq 0$, which is false. Thus,  $x$-values between -2 and 3 are \textit{not} in the domain of this function.
        \item For $x > 3$: plugging in $x = 4$ gives $(4)^2 - 4 -6 \geq 0$. When simplified, this gives $ 6 \geq 0$, which is true.
    \end{itemize}
    
    In interval notation, this is written as $$(-\infty, -2]\cup[3, \infty)$$

    \begin{quote} 
    \textit{Note: The brackets are used to show that the interval is inclusive of $-2$ and $3$ in the domain.}
    \end{quote}

\end{document}

